% ========================================================
% This document is a customizable CV/Resume template built using LaTeX.
% The template is designed for easy customization and clear structure.
%
% Author: Matthew DeVerna (www.matthewdeverna.com)
% Date: 2024
% Design: Cased on Hause Lin's CV (hauselin.com)
% 
% Project Overview:
% -----------------
% This LaTeX document is designed to help you create a professional CV or 
% resume with ease. It uses as little fancy LaTex functionality or custom functions as possible to maximize its longterm durability and flexibility.
% The document is structured into multiple sections, each loaded from 
% separate subfiles for modularity and ease of maintenance. 
%
% Key Features:
% -------------
% - Customizable sections: Education, Research Experience, Awards, Publications, etc.
% - Bookmarks in the PDF for easy navigation
% - Styled bibliography with BibLaTeX
% - Hyperlinked email and website
% - FontAwesome icons for additional styling
%
% Getting Started:
% ----------------
% 1. Customize your personal information by modifying the \mytitle command.
% 2. Add your content to the respective subfiles (e.g., education.tex, exp_research.tex).
% 3. Update the bibliography file (ref.bib) with your publications and categorize them with keywords.
%
% Important Notes:
% ----------------
% - This main file includes the overall structure and settings. 
% - Each section has its own detailed instructions for further customization.
%
% ========================================================

\documentclass[11pt]{article} % Choose the document class and font size
\usepackage[margin=1in]{geometry} % Page layout settings

% Set the citation style
\usepackage[
    backend=biber,     
    maxnames=20,        
    style=authoryear,       
    sorting=ydnt,  
    defernumbers=true,  
    uniquename=false,uniquelist=false
]{biblatex}
\addbibresource{ref.bib} % Adds the bibliography resource file 'ref.bib' containing all the references.


\usepackage{biblatex}
\DeclareFieldFormat{annotation}{\par\smallskip\textcolor{gray}{#1}\par}
\renewbibmacro*{annotation}{%
  \iffieldundef{annotation}
    {}
    {\printfield{annotation}}}
\usepackage{etoolbox}


% Patch the bibliography driver for 'article' type to add annotation after the entry
\AtEveryBibitem{%
  \ifentrytype{article}{%
    \csappto{blx@bbx@article}{\usebibmacro{annotation}}%
  }{}%
  \ifentrytype{inproceedings}{%
    \csappto{blx@bbx@inproceedings}{\usebibmacro{annotation}}%
  }{}%
  \ifentrytype{misc}{%
    \csappto{blx@bbx@misc}{\usebibmacro{annotation}}%
  }{}%
}

\DeclareFieldFormat[misc]{title}{``#1''}
\DeclareFieldFormat[misc]{howpublished}{\textit{#1}}
\DeclareFieldFormat[misc]{journaltitle}{\textit{#1}}
\renewbibmacro*{in:}{}

% remove alphabet after year
\AtEveryBibitem{%
  \ifnameundef{labelname}{}{%
    \clearfield{extradate}%
  }%
}



% ========================================================
% Allows columns that stretch across pages
\usepackage{longtable}
\newcommand{\rowspacing}{1.3} % default spacing multiplier
\renewcommand{\arraystretch}{1.3} % applies to all tables

% Before each longtable: negative vspace to pull it up
\preto\longtable{\vspace{-9pt}}

% After each longtable: small space
\appto\endlongtable{\vspace{-4pt}}

\usepackage{setspace} % for \setstretch
% Increase line spacing in bibliography
\AtBeginBibliography{\setstretch{1.3}} % 1.3× spacing

% Table functionality and beautification (not strictly needed)
\usepackage{bookmark}

% Use icons, if you want.
% All available icons: http://mirrors.ibiblio.org/CTAN/fonts/fontawesome5/doc/fontawesome5.pdf
\usepackage{fontawesome}

% Allows font justification control (needed for clean pub-list formatting)
\usepackage{ragged2e}

% For underlining with line breaks
\usepackage{soul} 

% All fonts: https://tug.org/FontCatalogue/
\usepackage{kpfonts} % More professional font
% \usepackage[default]{sourcecodepro} % Code-like font
\usepackage[T1]{fontenc}


% Control hyperlinks and colors
% CUSTOM COLORS INCLUDED DIRECTLY AFTER \begin{document}
\usepackage{xcolor}
\usepackage{hyperref}
\hypersetup{
    colorlinks=true,        % Enable colored links
    breaklinks=true,        % Allow links to break across lines
    linkcolor=cornflowerblue,    % Color of internal links
    urlcolor=cornflowerblue,     % Color of URL links
    anchorcolor=cornflowerblue,  % Color of anchors
    citecolor=cornflowerblue,    % Color of citations
    pdftitle={Mei-Yu Kuo - Curriculum Vitae},    % Title of the PDF
    pdfauthor={Mei-Yu Kuo}, % Author of the PDF
    bookmarksopen=true,      % Open bookmarks panel at start
}

%%% CONVENIENCE FUNCTIONS GO HERE %%%
%%% ----------------------------- %%%
\newcommand{\mytitle}[4]{
  \begin{center}
    \setstretch{1.1} % Adjust line spacing locally
    \Large\textbf{#1}\normalsize \\[2pt] % Name in large bold font
    \href{mailto:#2}{#2} \\[2pt] % Email with mailto link
    #3 \\[2pt] % Twitter/X or website
    #4 % Address
  \end{center}
}
% To hide years in Working papers % track the labels Under Reivew and In-prep
\AtEveryBibitem{%
  \ifboolexpr{
    test {\iffieldequalstr{howpublished}{Under Review}}
    or test {\iffieldequalstr{howpublished}{In preparation}}
    or test {\iffieldequalstr{howpublished}{R\&R}}
  }
  {%
    \clearfield{year}%
    \clearfield{date}%
    \clearfield{labelyear}%
    \renewbibmacro*{date}{}%
  }
  {}%
}
\AtEveryBibitem{%
  \iffieldequalstr{howpublished}{R\&R}{%
    \appto\finentry{%
      \smallskip
      \textit{ at \printfield{journaltitle}.}%
    }%
  }{}%
}

% Draw a horizontal rule under every \section*
\usepackage{titlesec}

% Define the horizontal rule
\newcommand{\sectionrule}{\vspace{2pt}\hrule\vspace{10pt}}

% Redefine unnumbered sections only
\titleformat{name=\section,numberless}
  {\large\bfseries}
  {}
  {0pt}
  {}
  [\sectionrule]


%%% ----------------------------- %%%


\begin{document}
% Set custom colors here (imported directly after \begin{document})
% The below use HTML hex codes.
% More HTML hex codes: https://encycolorpedia.com/html
\definecolor{firebrick}{HTML}{b22222} 
\definecolor{darkslategrey}{HTML}{2f4f4f} 
\definecolor{cornflowerblue}{HTML}{6495ed} 
\definecolor{mediumslateblue}{HTML}{7b68ee}  % Load custom colors from colors file
\mytitle{Mei-Yu Kuo}{kuo.355@buckeyemail.osu.edu}{\href{https://meisbase.github.io/}{Personal Website}}{1885 Neil Ave, Columbus, OH 43210} % Insert your custom title


% Ensure right side margin is not surpassed by bibliography and the right margin is aligned throughout
\RaggedRight


% These \pdfbookmark lines create bookmarks in the exported PDF document that display in the left pane.
% Value in [] sets the indentation level of the bookmark
\pdfbookmark[1]{Education}{}
\section*{Education}
% Add your educational background here!

% NOTE: If you want to remove the "Expected" footnote, you will want to remove:
% - Directly below: \renewcommand, \setcounter
% - In the table: \footnotemark in the left column
% - After the table: \footnotetext, \renewcommand, \setcounter

% Different numbers in "\setcounter{footnote}{0}" use different symbols
\renewcommand{\thefootnote}{\fnsymbol{footnote}}
\setcounter{footnote}{0}

\begin{longtable}[l]{@{}p{.125\textwidth} p{0.875\textwidth}}
    % Use custom symbol footnote for "expected"
    2023- & \textbf{Ph.D.}, Sociology, The Ohio State University  \\
          & \hspace{2em}- Graduate Specialization in Demography \\
          & \hspace{2em}- Affiliate, Institute for Population Research \\

    2025 & \textbf{M.A.}, Sociology, The Ohio State University \\
    2023 & \textbf{M.A.}, Sociology, National Taiwan University \\

    2021--22 & \textbf{Visiting Student}, University of Wisconsin-Madison \\

    2020 & \textbf{B.A.}, Journalism, National Chengchi University
% \end{tabularx}
\end{longtable}


% Add text for the custom footnote


% Restore the default footnote numbering
\renewcommand{\thefootnote}{\arabic{footnote}}
\setcounter{footnote}{1}



\pdfbookmark[1]{Research Interest}{interests}
\section*{Research Interests}
\label{interests}
Gender Inequality, Work \& Economy, Family, Migration

% \pdfbookmark[1]{Publications}{pubs}
% \section*{Research In-Progress}
% \label{pubs}

% Add equal contribution dagger

% \small
% \faGoogle~\href{https://scholar.google.com/}{Google Scholar}\\
% $\dagger \rightarrow$ Equal contribution
\normalsize


    % \pdfbookmark[2]{Journal Articles}{journal-article}
    % \subsection*{Journal Articles}
    % \label{journal-article}
    % \newrefcontext[labelprefix=J] % Will prefix bibliography numbers with this letter
    % % Ensures publications which are not cited in the document are included in the above sections
    % \nocite{*} 
    % \printbibliography[
    %     type=article, % Only include @article ref.bib items
    %     heading=none, % Do not include header. Gives us more control.
    %     resetnumbers=true, % Start item counter from zero
    %     keyword=J % Include items in ref.bib with keyword={J}
    % ]
    

   

    \pdfbookmark[2]{publications}{publications}
    \section*{Peer-Reviewed Publications}
    \label{article}
    \newrefcontext[sorting=none,labelprefix=R]
    \printbibliography[
      type=article,
      heading=none,
      resetnumbers=true,
      keyword=J
    ]

    \pdfbookmark[2]{bookchapters}{bookchapters}
   \section*{Book Chapters}
   \label{bookchapters}
   \newrefcontext[labelprefix=B]
   \nocite{*}
   \printbibliography[
      type=incollection,
      heading=none,
      resetnumbers=true,
      keyword=b
    ]
    
    \pdfbookmark[2]{workingpapers}{workingpapers}
    \section*{Working Papers}
    \textit{(* denotes equal authorship)}
    \label{workingpapers}
    \newrefcontext[sorting=none,labelprefix=R]
    \printbibliography[
      type=misc,
      heading=none,
      resetnumbers=true,
      keyword=R
    ]

    \pdfbookmark[2]{conference}{Conference}
    \section*{Conference Presentations}
    \label{conference}
    % Conference presentations, grouped by paper.
% Order: most recently presented paper first; within a paper, most recent venue first.
% One presentation:  \confsingle{authors}{title}{conference}{year}{format}
% Several:           \confpaper{authors}{title}  then one \venue{conference}{year}{format} each

\begingroup
\setstretch{1.3}
\newcommand{\confpaper}[2]{\par\vspace{4pt}\noindent\hangindent=1.5em\hangafter=1 #1. ``#2.''\par}
\newcommand{\venue}[3]{\par\noindent\hangindent=3.2em\hangafter=1\hspace*{1.5em}-- #1, #2, #3.\par}
\newcommand{\confsingle}[5]{\par\vspace{4pt}\noindent\hangindent=1.5em\hangafter=1 #1. ``#2.'' #3, #4, #5.\par}

\confpaper{Choi, Yoonyoung and \textbf{Mei-Yu Kuo}}{Parental Nativity Configuration and Adolescent Health: Disaggregating the Second Generation in the United States}
\venue{Annual Meeting of the American Sociological Association}{2026}{Oral Presentation}
\venue{Population Association of America Annual Meeting}{2026}{Flash Session}

\confpaper{\textbf{Kuo, Mei-Yu} and Vincent J. Roscigno}{Labor Force Attachment, College Major, and the Gender Earnings Gap}
\venue{The Society for the Advancement of Socio-Economics (SASE) Annual Conference}{2026}{Oral Presentation}
\venue{Population Association of America Annual Meeting}{2026}{Oral Presentation}
\venue{Annual Meeting of the American Sociological Association}{2025}{Roundtable Presentation}

\confsingle{\textbf{Kuo, Mei-Yu}, Vincent J. Roscigno, and Natasha Quadlin}{The Social-Psychological Costs of Workplace Gender Discrimination}{The Society for the Advancement of Socio-Economics (SASE) Annual Conference}{2026}{Oral Presentation}

\confsingle{\textbf{Kuo, Mei-Yu}}{Who Gains Equality? Education and Wives' Income Trajectories in Marriage}{Population Association of America Annual Meeting}{2026}{Flash Session}

\confsingle{\textbf{Kuo, Mei-Yu}, Vincent J. Roscigno, Ian Harwood, and Hui Zheng}{Workplace Age Discrimination and Its Implications for Mental Health}{Annual Meeting of the American Sociological Association}{2024}{Roundtable Presentation}

\confsingle{\textbf{Kuo, Mei-Yu}}{Delaying College Entry and Union Formation Patterns}{Population Association of America Annual Meeting}{2024}{Poster Presentation}

\confsingle{\textbf{Kuo, Mei-Yu}}{Women Candidates as a Contextual Factor on Women's Political Participation in Taiwan}{East Asian Junior Workshop}{2021}{Oral Presentation}
\endgroup
 % grouped by paper, written by hand

\pdfbookmark[1]{Awards \& Honors}{awards}
\section*{Awards \& Honors}
\label{awards}
% Include awards, honors, grant funding, etc. here

\begin{longtable}[l]{@{}p{.125\textwidth} p{0.875\textwidth}}
    % Use custom symbol footnote for "expected"

    2028-29  & University Fellowship for Dissertation Year, Graduate School, The Ohio State University (\$34,764) \\

    2026 &  Outstanding Second Year Paper Award, Department of sociology, The Ohio State University \\
    
    2025 & FREE Center Graduate Small Grant for Feminist Scholarship, Department of Women's, Gender and Sexuality Studies, The Ohio State University (\$500) \\
    
    2024 & FREE Center Graduate Small Grant for Feminist Scholarship, Department of Women's, Gender and Sexuality Studies, The Ohio State University (\$500) \\
    
    2024 & Best Thesis Award, Population Association of Taiwan \\
    
    2024 & Summer Methods Award for Method Training, Department of Sociology, The Ohio State University (\$1,500) \\
    
    2024 & Poster Award, Family Demography Session, 2024 Population Association of America Annual Meeting \\
    
    2023--24 & University Fellowship, Graduate School, The Ohio State University (\$30,420) \\
    
    2021--22 & Study Abroad Scholarship for Future Scholars in the Humanities and Social Sciences, Ministry of Education, Taiwan  (\$31,000) \\

    2018--20 & Undergraduate Presidential Award * 4, Department of Journalism, National Chengchi University \\

    2020 & Pan Sih-Yuan Funding for International Reporting Project, Department of Journalism, National Chengchi University (\$3,300) 


    % Add more things here!
   
\end{longtable}


% Research Experience section hidden (content kept in exp_research.tex)
% \pdfbookmark[1]{Research Experience}{exp_research}
% \section*{Research Experience}
% \label{exp_research}
% % List your research experience here

\begin{longtable}[l]{@{}p{.125\textwidth} p{0.875\textwidth}}

    2025 & Research Associate, The Ohio State University (Supervisor: \href{https://www.samclark.net/}{Dr. Samuel J. Clark}) \\
    
    2025 & Research Associate, The Ohio State University (Supervisor: \href{https://u.osu.edu/roscigno.1/}{Dr. Vincent J. Roscigno} \& \href{https://www.huizhenghome.com/}{Dr. Hui Zheng}) \\

    2024 & Research Associate, The Ohio State University (Supervisor: \href{https://www.huizhenghome.com/}{Dr. Hui Zheng}) \\

    2023 & Research Associate, Academia Sinica (Supervisor: \href{https://sites.google.com/view/hsuanweilee/home}{Dr. Hsuan-Wei "Wayne" Lee}) \\

\end{longtable}

\pdfbookmark[1]{Teaching}{teaching}
\section*{Teaching}
\label{teaching}


\begin{longtable}[l]{@{}p{.125\textwidth} p{0.875\textwidth}}

    2026 & Teaching Associate, Research Methods in Sociology, The Ohio State University (Supervisor: \href{http://www.fangqiwen.com/}{Dr. Fangqi Wen}) \\
    
    2024 & Teaching Associate, Statistics in Sociology, The Ohio State University (Supervisor: \href{https://u.osu.edu/melamed9/}{Dr. David Melamed}) \\

    2023 & Lecturer, Research Methods in Social Science (1-hour session), High School Humanities \& Social Sciences Summer Program, Ministry of Education, Taiwan \\

    2021 & Teaching Associate, Statistics in Sociology, National Taiwan University (Supervisor: \href{https://scholar.google.com/citations?hl=en&user=Lpco_MgAAAAJ&view_op=list_works&sortby=pubdate}{Dr. Janet Chen-Lan Kuo}) \\


\end{longtable}


\pdfbookmark[1]{Academic Service}{service}
\section*{Professional Service}
\label{service}


% \subsection*{Guest Editor}
% 
\begin{longtable}[l]{@{}p{.125\textwidth} p{0.875\textwidth}}


    2025--26 & President, Institute for Population Research Graduate Student Organization, The Ohio State University \\

    2024--25 & Secretary, Institute for Population Research Graduate Student Organization, The Ohio State University \\

% \end{tabularx}
\end{longtable}


% \subsection*{Journal Reviewer}
% % 
% \begin{longtable}[l]{@{}p{.125\textwidth} p{0.875\textwidth}}

%     2023 & Nature\\
%     2023 & Science \\
    
% % \end{tabularx}
% \end{longtable}




\pdfbookmark[1]{Professional Affiliation}{affiliation}
\section*{Professional Affiliation}
\label{affiliation}
\begin{longtable}{@{}p{\textwidth}@{}}
American Sociological Association; \\[2pt]
Population Association of America; \\[2pt]
The Society for the Advancement of Socio-Economics; \\[2pt]
\href{https://webpageprodvm.ntu.edu.tw/soc.quant.lab/Default.aspx}{Quantitative Research Lab}, Department of Sociology, National Taiwan University
\end{longtable}


\pdfbookmark[1]{Tools \& Software}{tools}
\section*{Skills}
\label{tools}
\begin{spacing}{1.3}
\textbf{Language:} English (Fluent), Mandarin Chinese (Native) \\
\textbf{Computing:} R, STATA, Python, \LaTeX \\
\end{spacing}




% Include any additional details here
% \vspace{-.75em}
% \small
% $\dagger \rightarrow$ Equal contribution
% \normalsize

% \pdfbookmark[2]{Talks}{talks}
% \subsection*{Talks}
% \label{talks}
% \newrefcontext[labelprefix=T]
% \printbibliography[type=misc,heading=none,resetnumbers=true,keyword=T]

% \pdfbookmark[2]{Posters}{posters}
% \subsection*{Posters}
% \label{posters}
% \newrefcontext[labelprefix=P]        
% \printbibliography[type=misc,heading=none,resetnumbers=true,keyword=P]

% \pdfbookmark[2]{Demonstrations \& Tutorials}{demos}
% \subsection*{Demonstrations \& Tutorials}
% \label{demos}
% \newrefcontext[labelprefix=D] \printbibliography[type=misc,heading=none,resetnumbers=true,keyword=D]


% \pdfbookmark[1]{Selected Media Coverage}{media}
% \section*{Selected Media Coverage}
% \label{media}
% %  Include media coverage here
% 
\begin{longtable}[l]{@{}p{.125\textwidth} p{0.875\textwidth}}

    2024 & The New York Times, \href{https://nyt.com}{Scientists figure out the meaning of life.}~\cite{art1} (Note how the regular citation functions work throughout the document!) \\
    
    2024 & AP News, \href{https://apnews.com/}{What even is science?}~\cite{art2} (Point to any project in your ref.bib file with the standard `cite' function) \\

    
\end{longtable}

 % Input lines load the material from the subdocuments


% \pdfbookmark[1]{Teaching}{teaching}
% \section*{Teaching}
% \label{teaching}
% 

\begin{longtable}[l]{@{}p{.125\textwidth} p{0.875\textwidth}}

    2026 & Teaching Associate, Research Methods in Sociology, The Ohio State University (Supervisor: \href{http://www.fangqiwen.com/}{Dr. Fangqi Wen}) \\
    
    2024 & Teaching Associate, Statistics in Sociology, The Ohio State University (Supervisor: \href{https://u.osu.edu/melamed9/}{Dr. David Melamed}) \\

    2023 & Lecturer, Research Methods in Social Science (1-hour session), High School Humanities \& Social Sciences Summer Program, Ministry of Education, Taiwan \\

    2021 & Teaching Associate, Statistics in Sociology, National Taiwan University (Supervisor: \href{https://scholar.google.com/citations?hl=en&user=Lpco_MgAAAAJ&view_op=list_works&sortby=pubdate}{Dr. Janet Chen-Lan Kuo}) \\


\end{longtable}


% \pdfbookmark[1]{Academic Advising}{advising}
% \section*{Academic Advising}
% \label{advising}
% % Including academic advisement history here


% Remove \subsection{} lines and multiple tables if you only need one section!
\subsection*{Graduate}
\begin{longtable}[l]{@{}p{.125\textwidth} p{0.875\textwidth}}

    2023 & Name, University, Additional details \\

\end{longtable}

\subsection*{Undergraduate}
\begin{longtable}[l]{@{}p{.125\textwidth} p{0.875\textwidth}}

    2023 & Name, University, Additional details \\

\end{longtable}


% \pdfbookmark[1]{Other Experience}{exp_other}
% \section*{Other Experience}
% \label{exp_other}
% % Add other experience you don't know where to classify here

\begin{longtable}[l]{@{}p{.125\textwidth} p{0.875\textwidth}}

    2018 & Past Non-academic Job, \href{https://example.com/}{Company Name}, New York City, New York (Additional detail) \\

\end{longtable}


% Pretty ending with the date last updated
\centering
\rule{0.25\linewidth}{0.4pt}\\
\medskip
Last updated: Spring 2026

\end{document}
